\section{RISC-V CoVE / AP-TEE confidential VM}

\begin{frame}[t,shrink=6]{RISC-V CoVE / AP-TEE confidential VM：方向开场}
\DeckKicker{方向开场}
\SlideMeta{证据等级}{方向综述}{来源状态：公开材料综合}
{\large\bfseries\color{Ink} RISC-V CoVE / AP-TEE confidential VM 的核心是先界定保护对象、管理边界和证据强度，再用三篇主讲材料串起技术演进。\par}\vspace{1.8mm}
\begin{columns}[T,totalwidth=\textwidth]
  \begin{column}{0.45\textwidth}
    \fcolorbox{Rule}{Paper}{%
  \begin{minipage}[t]{0.97\linewidth}
    \vspace{1.1mm}
    \hspace{1.4mm}\begin{minipage}{0.92\linewidth}
      \PanelTitle{内容说明}
      \scriptsize \begin{itemize}
  \item 把 pre-CoVE enclave 与 AP-TEE/CoVE TVM lifecycle 区分开。
  \item 固定三篇主讲材料；survey 只作为 background substrate，不替代 CoVE/AP-TEE 机制证据。
  \item 先讲方向痛点和设计轴，再逐篇说明贡献、限制、证据强度和可落地条件。
  \item 对规范、Survey、draft、vendor evidence 保持显式标注，避免把背景材料写成一手系统证明。
\end{itemize}
    \end{minipage}
    \vspace{1mm}
  \end{minipage}}
  \end{column}
  \begin{column}{0.49\textwidth}
    \fcolorbox{Rule}{Panel}{%
  \begin{minipage}[t]{0.97\linewidth}
    \vspace{1.1mm}
    \hspace{1.4mm}\begin{minipage}{0.92\linewidth}
      \PanelTitle{三篇主讲材料的分工}
      \scriptsize {\tiny
\begin{tabularx}{\linewidth}{@{}YYY@{}}
\toprule
\textbf{论文} & \textbf{定位} & \textbf{证据边界} \\
\midrule
RISC-V CoVE reference architecture & 开创/基础入口 & E3 / 基础证据 public preprint \\
RISC-V AP-TEE / CoVE TVM 规范草案 & 代表性改进 & E3 / 草案/未批准规范 \\
RISC-V TEE 谱系辅助材料 & 当前 SOTA/补充边界 & E2 survey / Background substrate \\
\bottomrule
\end{tabularx}
}
    \end{minipage}
    \vspace{1mm}
  \end{minipage}}
  \end{column}
\end{columns}
\vspace{1mm}
{\tiny\textcolor{Muted}{引用依据：本页综合三篇主讲材料的研究定位、证据等级和公开来源状态。}}
\end{frame}


\begin{frame}[t,shrink=6]{RISC-V CoVE reference architecture：内容摘要}
\DeckKicker{内容摘要}
\SlideMeta{证据等级}{E3 / 基础证据 public preprint}{来源状态：本地 PDF 已核验}
{\large\bfseries\color{Ink} 提出 RISC-V CoVE confidential computing 参考架构方向，明确 confidential VM 需要的 ISA、n。\par}\vspace{1.8mm}
\begin{columns}[T,totalwidth=\textwidth]
  \begin{column}{0.45\textwidth}
    \fcolorbox{Rule}{Paper}{%
  \begin{minipage}[t]{0.97\linewidth}
    \vspace{1.1mm}
    \hspace{1.4mm}\begin{minipage}{0.92\linewidth}
      \PanelTitle{内容说明}
      \scriptsize \begin{itemize}
  \item 提出 RISC-V CoVE confidential computing 参考架构方向，明确 confidential VM 需要的 ISA、non-ISA、SoC 与 platf。
  \item CoVE 2023 is the early mainline design material for RISC-V confidential VMs.
\end{itemize}
    \end{minipage}
    \vspace{1mm}
  \end{minipage}}
  \end{column}
  \begin{column}{0.49\textwidth}
    \fcolorbox{Rule}{Panel}{%
  \begin{minipage}[t]{0.97\linewidth}
    \vspace{1.1mm}
    \hspace{1.4mm}\begin{minipage}{0.92\linewidth}
      \PanelTitle{证据定位}
      \scriptsize {\tiny
\begin{tabularx}{\linewidth}{@{}YY@{}}
\toprule
\textbf{字段} & \textbf{内容} \\
\midrule
论文定位 & 开创/基础入口 \\
材料类型 & 系统论文 \\
证据等级 & E3 / 基础证据 public preprint \\
来源状态 & 本地 PDF 已核验 \\
\bottomrule
\end{tabularx}
}
    \end{minipage}
    \vspace{1mm}
  \end{minipage}}
  \end{column}
\end{columns}
\vspace{1mm}
{\tiny\textcolor{Muted}{引用依据：引用依据为论文原文、官方规范或公开材料；本页结论按 E3 / 基础证据 public preprint 的证据等级使用，不外推到未验证平台。}}
\end{frame}


\begin{frame}[t,shrink=6]{RISC-V CoVE reference architecture：研究背景}
\DeckKicker{研究背景}
\SlideMeta{证据等级}{E3 / 基础证据 public preprint}{来源状态：本地 PDF 已核验}
{\large\bfseries\color{Ink} 多租户 confidential VM 需要降低 tenant 对 host/hypervisor 的信任，并区别于单进程 enclave lin。\par}\vspace{1.8mm}
\begin{columns}[T,totalwidth=\textwidth]
  \begin{column}{0.45\textwidth}
    \fcolorbox{Rule}{Paper}{%
  \begin{minipage}[t]{0.97\linewidth}
    \vspace{1.1mm}
    \hspace{1.4mm}\begin{minipage}{0.92\linewidth}
      \PanelTitle{内容说明}
      \scriptsize \begin{itemize}
  \item 多租户 confidential VM 需要降低 tenant 对 host/hypervisor 的信任，并区别于单进程 enclave lineage
\end{itemize}
    \end{minipage}
    \vspace{1mm}
  \end{minipage}}
  \end{column}
  \begin{column}{0.49\textwidth}
    \fcolorbox{Rule}{Panel}{%
  \begin{minipage}[t]{0.97\linewidth}
    \vspace{1.1mm}
    \hspace{1.4mm}\begin{minipage}{0.92\linewidth}
      \PanelTitle{问题边界}
      \scriptsize {\tiny
\begin{tabularx}{\linewidth}{@{}YY@{}}
\toprule
\textbf{维度} & \textbf{说明} \\
\midrule
保护对象 & RISC-V CoVE / AP-TEE confidential VM \\
传统瓶颈 & 把 pre-CoVE enclave 与 AP-TEE/CoVE TVM lifecycle 区分开。 \\
本文入口 & RISC-V CoVE reference architecture \\
证据限制 & E3 / 基础证据 public preprint \\
\bottomrule
\end{tabularx}
}
    \end{minipage}
    \vspace{1mm}
  \end{minipage}}
  \end{column}
\end{columns}
\vspace{1mm}
{\tiny\textcolor{Muted}{引用依据：引用依据为论文原文、官方规范或公开材料；本页结论按 E3 / 基础证据 public preprint 的证据等级使用，不外推到未验证平台。}}
\end{frame}


\begin{frame}[t,shrink=6]{RISC-V CoVE reference architecture：解决方案}
\DeckKicker{解决方案}
\SlideMeta{证据等级}{E3 / 基础证据 public preprint}{来源状态：本地 PDF 已核验}
{\large\bfseries\color{Ink} 用 CoVE reference architecture 描述 TVM 所需的硬件、固件、hypervisor 和 platform suppo。\par}\vspace{1.8mm}
\begin{columns}[T,totalwidth=\textwidth]
  \begin{column}{0.45\textwidth}
    \fcolorbox{Rule}{Paper}{%
  \begin{minipage}[t]{0.97\linewidth}
    \vspace{1.1mm}
    \hspace{1.4mm}\begin{minipage}{0.92\linewidth}
      \PanelTitle{内容说明}
      \scriptsize \begin{itemize}
  \item 用 CoVE reference architecture 描述 TVM 所需的硬件、固件、hypervisor 和 platform support
\end{itemize}
    \end{minipage}
    \vspace{1mm}
  \end{minipage}}
  \end{column}
  \begin{column}{0.49\textwidth}
    \fcolorbox{Rule}{Panel}{%
  \begin{minipage}[t]{0.97\linewidth}
    \vspace{1.1mm}
    \hspace{1.4mm}\begin{minipage}{0.92\linewidth}
      \PanelTitle{核心设计拆解}
      \scriptsize \begin{itemize}
  \item 01. RISC-V confidential VM reference architecture
  \item 02. ISA/non-ISA/SoC/platform requirement split
  \item 03. Host/hypervisor trust reduction for TVM
\end{itemize}
    \end{minipage}
    \vspace{1mm}
  \end{minipage}}
  \end{column}
\end{columns}
\vspace{1mm}
{\tiny\textcolor{Muted}{引用依据：引用依据为论文原文、官方规范或公开材料；本页结论按 E3 / 基础证据 public preprint 的证据等级使用，不外推到未验证平台。}}
\end{frame}


\begin{frame}[t,shrink=6]{RISC-V CoVE reference architecture：实验结果}
\DeckKicker{实验结果}
\SlideMeta{证据等级}{E3 / 基础证据 public preprint}{来源状态：本地 PDF 已核验}
{\large\bfseries\color{Ink} 架构论文/position-style，无完整系统实验\par}\vspace{1.8mm}
\begin{columns}[T,totalwidth=\textwidth]
  \begin{column}{0.45\textwidth}
    \fcolorbox{Rule}{Paper}{%
  \begin{minipage}[t]{0.97\linewidth}
    \vspace{1.1mm}
    \hspace{1.4mm}\begin{minipage}{0.92\linewidth}
      \PanelTitle{内容说明}
      \scriptsize \begin{itemize}
  \item 架构论文/position-style，无完整系统实验
  \item 具体 ABI/state details 需要 AP-TEE draft
\end{itemize}
    \end{minipage}
    \vspace{1mm}
  \end{minipage}}
  \end{column}
  \begin{column}{0.49\textwidth}
    \fcolorbox{Rule}{Panel}{%
  \begin{minipage}[t]{0.97\linewidth}
    \vspace{1.1mm}
    \hspace{1.4mm}\begin{minipage}{0.92\linewidth}
      \PanelTitle{实验/证据边界}
      \scriptsize {\tiny
\begin{tabularx}{\linewidth}{@{}YYY@{}}
\toprule
\textbf{对象} & \textbf{可支撑} & \textbf{边界} \\
\midrule
数据/来源 & RISC-V CoVE reference architecture & 本地 PDF 已核验 \\
Baseline/对照 & 以论文或规范原文为准 & 不跨材料外推 \\
指标/对象 & 只使用当前来源可证明的信息 & E3 / 基础证据 public preprint \\
使用边界 & 可进入本方向演进图 & 不能替代更强证据 \\
\bottomrule
\end{tabularx}
}
    \end{minipage}
    \vspace{1mm}
  \end{minipage}}
  \end{column}
\end{columns}
\vspace{1mm}
{\tiny\textcolor{Muted}{引用依据：引用依据为论文原文、官方规范或公开材料；本页结论按 E3 / 基础证据 public preprint 的证据等级使用，不外推到未验证平台。}}
\end{frame}


\begin{frame}[t,shrink=6]{RISC-V CoVE reference architecture：文章评价}
\DeckKicker{文章评价}
\SlideMeta{证据等级}{E3 / 基础证据 public preprint}{来源状态：本地 PDF 已核验}
{\large\bfseries\color{Ink} CoVE confidential VM 入口材料，能把 RISC-V 从 enclave lineage 推向 TVM 模型\par}\vspace{1.8mm}
\begin{columns}[T,totalwidth=\textwidth]
  \begin{column}{0.45\textwidth}
    \fcolorbox{Rule}{Paper}{%
  \begin{minipage}[t]{0.97\linewidth}
    \vspace{1.1mm}
    \hspace{1.4mm}\begin{minipage}{0.92\linewidth}
      \PanelTitle{内容说明}
      \scriptsize \begin{itemize}
  \item 设计优势：CoVE confidential VM 入口材料，能把 RISC-V 从 enclave lineage 推向 TVM 模型。
  \item 局限性：不提供完整 ABI、状态机或 ratified standard 语义。
  \item 商业化潜力：对开放 confidential VM 生态有方向价值；风险在标准未定、平台支持和 verifier ecosystem。
  \item 证据边界：E3 / Foundational public preprint；本地 PDF 已核验。
\end{itemize}
    \end{minipage}
    \vspace{1mm}
  \end{minipage}}
  \end{column}
  \begin{column}{0.49\textwidth}
    \fcolorbox{Rule}{Panel}{%
  \begin{minipage}[t]{0.97\linewidth}
    \vspace{1.1mm}
    \hspace{1.4mm}\begin{minipage}{0.92\linewidth}
      \PanelTitle{文章评价}
      \scriptsize {\tiny
\begin{tabularx}{\linewidth}{@{}YY@{}}
\toprule
\textbf{维度} & \textbf{判断} \\
\midrule
设计优势 & CoVE confidential VM 入口材料，能把 RISC-V 从 enclave lineage 推向 TVM 模型。 \\
主要局限 & 不提供完整 ABI、状态机或 ratified standard 语义。 \\
商业落地 & 对开放 confidential VM 生态有方向价值；风险在标准未定、平台支持和 verifier ecosystem。 \\
讲稿定位 & 开创/基础入口; 证据强度 E3 \\
\bottomrule
\end{tabularx}
}
    \end{minipage}
    \vspace{1mm}
  \end{minipage}}
  \end{column}
\end{columns}
\vspace{1mm}
{\tiny\textcolor{Muted}{引用依据：引用依据为论文原文、官方规范或公开材料；本页结论按 E3 / 基础证据 public preprint 的证据等级使用，不外推到未验证平台。}}
\end{frame}


\begin{frame}[t,shrink=6]{RISC-V AP-TEE / CoVE TVM 规范草案：内容摘要}
\DeckKicker{内容摘要}
\SlideMeta{证据等级}{E3 / 草案/未批准规范}{来源状态：本地 PDF 已核验}
{\large\bfseries\color{Ink} 定义 RISC-V AP-TEE/CoVE TVM lifecycle 和 SBI ABI，包括 TVM、TSM、COVH/COVG 与 memo。\par}\vspace{1.8mm}
\begin{columns}[T,totalwidth=\textwidth]
  \begin{column}{0.45\textwidth}
    \fcolorbox{Rule}{Paper}{%
  \begin{minipage}[t]{0.97\linewidth}
    \vspace{1.1mm}
    \hspace{1.4mm}\begin{minipage}{0.92\linewidth}
      \PanelTitle{内容说明}
      \scriptsize \begin{itemize}
  \item 定义 RISC-V AP-TEE/CoVE TVM lifecycle 和 SBI ABI，包括 TVM、TSM、COVH/COVG 与 memory lifecycle
  \item AP-TEE draft is the current standards-track source for RISC-V confidential VM semantics.
\end{itemize}
    \end{minipage}
    \vspace{1mm}
  \end{minipage}}
  \end{column}
  \begin{column}{0.49\textwidth}
    \fcolorbox{Rule}{Panel}{%
  \begin{minipage}[t]{0.97\linewidth}
    \vspace{1.1mm}
    \hspace{1.4mm}\begin{minipage}{0.92\linewidth}
      \PanelTitle{证据定位}
      \scriptsize {\tiny
\begin{tabularx}{\linewidth}{@{}YY@{}}
\toprule
\textbf{字段} & \textbf{内容} \\
\midrule
论文定位 & 代表性改进 \\
材料类型 & 规范/标准材料 \\
证据等级 & E3 / 草案/未批准规范 \\
来源状态 & 本地 PDF 已核验 \\
\bottomrule
\end{tabularx}
}
    \end{minipage}
    \vspace{1mm}
  \end{minipage}}
  \end{column}
\end{columns}
\vspace{1mm}
{\tiny\textcolor{Muted}{引用依据：引用依据为论文原文、官方规范或公开材料；本页结论按 E3 / 草案/未批准规范 的证据等级使用，不外推到未验证平台。}}
\end{frame}


\begin{frame}[t,shrink=6]{RISC-V AP-TEE / CoVE TVM 规范草案：研究背景}
\DeckKicker{研究背景}
\SlideMeta{证据等级}{E3 / 草案/未批准规范}{来源状态：本地 PDF 已核验}
{\large\bfseries\color{Ink} CoVE 需要标准化 TVM、TSM、memory lifecycle、attestation 和 host/guest 调用界面，才能与 Arm。\par}\vspace{1.8mm}
\begin{columns}[T,totalwidth=\textwidth]
  \begin{column}{0.45\textwidth}
    \fcolorbox{Rule}{Paper}{%
  \begin{minipage}[t]{0.97\linewidth}
    \vspace{1.1mm}
    \hspace{1.4mm}\begin{minipage}{0.92\linewidth}
      \PanelTitle{内容说明}
      \scriptsize \begin{itemize}
  \item CoVE 需要标准化 TVM、TSM、memory lifecycle、attestation 和 host/guest 调用界面，才能与 Arm CCA/RMM 同层比较
\end{itemize}
    \end{minipage}
    \vspace{1mm}
  \end{minipage}}
  \end{column}
  \begin{column}{0.49\textwidth}
    \fcolorbox{Rule}{Panel}{%
  \begin{minipage}[t]{0.97\linewidth}
    \vspace{1.1mm}
    \hspace{1.4mm}\begin{minipage}{0.92\linewidth}
      \PanelTitle{问题边界}
      \scriptsize {\tiny
\begin{tabularx}{\linewidth}{@{}YY@{}}
\toprule
\textbf{维度} & \textbf{说明} \\
\midrule
保护对象 & RISC-V CoVE / AP-TEE confidential VM \\
传统瓶颈 & 把 pre-CoVE enclave 与 AP-TEE/CoVE TVM lifecycle 区分开。 \\
本文入口 & RISC-V AP-TEE / CoVE TVM 规范草案 \\
证据限制 & E3 / 草案/未批准规范 \\
\bottomrule
\end{tabularx}
}
    \end{minipage}
    \vspace{1mm}
  \end{minipage}}
  \end{column}
\end{columns}
\vspace{1mm}
{\tiny\textcolor{Muted}{引用依据：引用依据为论文原文、官方规范或公开材料；本页结论按 E3 / 草案/未批准规范 的证据等级使用，不外推到未验证平台。}}
\end{frame}


\begin{frame}[t,shrink=6]{RISC-V AP-TEE / CoVE TVM 规范草案：解决方案}
\DeckKicker{解决方案}
\SlideMeta{证据等级}{E3 / 草案/未批准规范}{来源状态：本地 PDF 已核验}
{\large\bfseries\color{Ink} 给出 TSM、TSM-driver、Supervisor Domains、COVH/COVG、memory donation/re结论/sh。\par}\vspace{1.8mm}
\begin{columns}[T,totalwidth=\textwidth]
  \begin{column}{0.45\textwidth}
    \fcolorbox{Rule}{Paper}{%
  \begin{minipage}[t]{0.97\linewidth}
    \vspace{1.1mm}
    \hspace{1.4mm}\begin{minipage}{0.92\linewidth}
      \PanelTitle{内容说明}
      \scriptsize \begin{itemize}
  \item 给出 TSM、TSM-driver、Supervisor Domains、COVH/COVG、memory donation/re结论/share 等 draft 规范语义
\end{itemize}
    \end{minipage}
    \vspace{1mm}
  \end{minipage}}
  \end{column}
  \begin{column}{0.49\textwidth}
    \fcolorbox{Rule}{Panel}{%
  \begin{minipage}[t]{0.97\linewidth}
    \vspace{1.1mm}
    \hspace{1.4mm}\begin{minipage}{0.92\linewidth}
      \PanelTitle{核心设计拆解}
      \scriptsize \begin{itemize}
  \item 01. TVM/TSM lifecycle and SBI ABI
  \item 02. COVH/COVG host/guest interface split
  \item 03. Memory donation, re结论 and sharing semantics
\end{itemize}
    \end{minipage}
    \vspace{1mm}
  \end{minipage}}
  \end{column}
\end{columns}
\vspace{1mm}
{\tiny\textcolor{Muted}{引用依据：引用依据为论文原文、官方规范或公开材料；本页结论按 E3 / 草案/未批准规范 的证据等级使用，不外推到未验证平台。}}
\end{frame}


\begin{frame}[t,shrink=6]{RISC-V AP-TEE / CoVE TVM 规范草案：实验结果}
\DeckKicker{实验结果}
\SlideMeta{证据等级}{E3 / 草案/未批准规范}{来源状态：本地 PDF 已核验}
{\large\bfseries\color{Ink} 规范/Survey，无新实验；RISC-V AP-TEE / CoVE TVM 规范草案 只支撑术语、taxonomy、接口语义或证据边界。\par}\vspace{1.8mm}
\begin{columns}[T,totalwidth=\textwidth]
  \begin{column}{0.45\textwidth}
    \fcolorbox{Rule}{Paper}{%
  \begin{minipage}[t]{0.97\linewidth}
    \vspace{1.1mm}
    \hspace{1.4mm}\begin{minipage}{0.92\linewidth}
      \PanelTitle{内容说明}
      \scriptsize \begin{itemize}
  \item 本页不展示独立系统实验，也不把二手归纳写成一手结果。
  \item 可支撑：AP-TEE draft is the current standards-track source for RISC-V confidential VM semanti。
  \item 证据等级：E3 / Draft-not-ratified；来源状态：本地 PDF 已核验。
  \item 涉及具体平台、协议状态或数值时，转到原始论文、官方规范或厂商材料核验。
\end{itemize}
    \end{minipage}
    \vspace{1mm}
  \end{minipage}}
  \end{column}
  \begin{column}{0.49\textwidth}
    \fcolorbox{Rule}{Panel}{%
  \begin{minipage}[t]{0.97\linewidth}
    \vspace{1.1mm}
    \hspace{1.4mm}\begin{minipage}{0.92\linewidth}
      \PanelTitle{实验/证据边界}
      \scriptsize {\tiny
\begin{tabularx}{\linewidth}{@{}YYY@{}}
\toprule
\textbf{对象} & \textbf{可支撑} & \textbf{边界} \\
\midrule
数据/来源 & RISC-V AP-TEE / CoVE TVM 规范草案 & 本地 PDF 已核验 \\
Baseline/对照 & 以论文或规范原文为准 & 不跨材料外推 \\
指标/对象 & 只使用当前来源可证明的信息 & E3 / 草案/未批准规范 \\
使用边界 & 可进入本方向演进图 & 不能替代更强证据 \\
\bottomrule
\end{tabularx}
}
    \end{minipage}
    \vspace{1mm}
  \end{minipage}}
  \end{column}
\end{columns}
\vspace{1mm}
{\tiny\textcolor{Muted}{引用依据：引用依据为论文原文、官方规范或公开材料；本页结论按 E3 / 草案/未批准规范 的证据等级使用，不外推到未验证平台。}}
\end{frame}


\begin{frame}[t,shrink=6]{RISC-V AP-TEE / CoVE TVM 规范草案：文章评价}
\DeckKicker{文章评价}
\SlideMeta{证据等级}{E3 / 草案/未批准规范}{来源状态：本地 PDF 已核验}
{\large\bfseries\color{Ink} 评价：RISC-V AP-TEE / CoVE TVM 规范草案 适合做 taxonomy/规范入口；规范/Survey，无新实验，不能替代一手系统证据。\par}\vspace{1.8mm}
\begin{columns}[T,totalwidth=\textwidth]
  \begin{column}{0.45\textwidth}
    \fcolorbox{Rule}{Paper}{%
  \begin{minipage}[t]{0.97\linewidth}
    \vspace{1.1mm}
    \hspace{1.4mm}\begin{minipage}{0.92\linewidth}
      \PanelTitle{内容说明}
      \scriptsize \begin{itemize}
  \item 设计优势：当前最关键的 RISC-V CCA 对照材料，直接定义 TVM/TSM lifecycle。
  \item 局限性：v0.7 RC2 仍是 draft，状态机和 ABI 可能变化。
  \item 商业化潜力：商业潜力在开放 confidential VM stack；风险在标准收敛、firmware/OS support 和兼容迁移。
  \item 规范/Survey，无新实验；具体平台结论转到一手材料核验。
\end{itemize}
    \end{minipage}
    \vspace{1mm}
  \end{minipage}}
  \end{column}
  \begin{column}{0.49\textwidth}
    \fcolorbox{Rule}{Panel}{%
  \begin{minipage}[t]{0.97\linewidth}
    \vspace{1.1mm}
    \hspace{1.4mm}\begin{minipage}{0.92\linewidth}
      \PanelTitle{文章评价}
      \scriptsize {\tiny
\begin{tabularx}{\linewidth}{@{}YY@{}}
\toprule
\textbf{维度} & \textbf{判断} \\
\midrule
设计优势 & 当前最关键的 RISC-V CCA 对照材料，直接定义 TVM/TSM lifecycle。 \\
主要局限 & v0.7 RC2 仍是 draft，状态机和 ABI 可能变化。 \\
商业落地 & 商业潜力在开放 confidential VM stack；风险在标准收敛、firmware/OS support 和兼容迁移。 \\
讲稿定位 & 代表性改进; 证据强度 E3 \\
\bottomrule
\end{tabularx}
}
    \end{minipage}
    \vspace{1mm}
  \end{minipage}}
  \end{column}
\end{columns}
\vspace{1mm}
{\tiny\textcolor{Muted}{引用依据：引用依据为论文原文、官方规范或公开材料；本页结论按 E3 / 草案/未批准规范 的证据等级使用，不外推到未验证平台。}}
\end{frame}


\begin{frame}[t,shrink=6]{RISC-V TEE 谱系辅助材料：内容摘要}
\DeckKicker{内容摘要}
\SlideMeta{证据等级}{E2 survey / Background substrate}{来源状态：本地 PDF 已核验}
{\large\bfseries\color{Ink} 把 pre-CoVE enclave 与 CoVE 放入 RISC-V TEE 谱系，辅助区分 enclave architecture、trus。\par}\vspace{1.8mm}
\begin{columns}[T,totalwidth=\textwidth]
  \begin{column}{0.45\textwidth}
    \fcolorbox{Rule}{Paper}{%
  \begin{minipage}[t]{0.97\linewidth}
    \vspace{1.1mm}
    \hspace{1.4mm}\begin{minipage}{0.92\linewidth}
      \PanelTitle{内容说明}
      \scriptsize \begin{itemize}
  \item 把 pre-CoVE enclave 与 CoVE 放入 RISC-V TEE 谱系，辅助区分 enclave architecture、trusted monitor/runtim。
  \item Recent RISC-V TEE survey cross-checks the CoVE and AP-TEE lineage against older enclave sys。
\end{itemize}
    \end{minipage}
    \vspace{1mm}
  \end{minipage}}
  \end{column}
  \begin{column}{0.49\textwidth}
    \fcolorbox{Rule}{Panel}{%
  \begin{minipage}[t]{0.97\linewidth}
    \vspace{1.1mm}
    \hspace{1.4mm}\begin{minipage}{0.92\linewidth}
      \PanelTitle{证据定位}
      \scriptsize {\tiny
\begin{tabularx}{\linewidth}{@{}YY@{}}
\toprule
\textbf{字段} & \textbf{内容} \\
\midrule
论文定位 & 当前 SOTA/补充边界 \\
材料类型 & SoK/综述 \\
证据等级 & E2 survey / Background substrate \\
来源状态 & 本地 PDF 已核验 \\
\bottomrule
\end{tabularx}
}
    \end{minipage}
    \vspace{1mm}
  \end{minipage}}
  \end{column}
\end{columns}
\vspace{1mm}
{\tiny\textcolor{Muted}{引用依据：引用依据为论文原文、官方规范或公开材料；本页结论按 E2 survey / Background substrate 的证据等级使用，不外推到未验证平台。}}
\end{frame}


\begin{frame}[t,shrink=6]{RISC-V TEE 谱系辅助材料：研究背景}
\DeckKicker{研究背景}
\SlideMeta{证据等级}{E2 survey / Background substrate}{来源状态：本地 PDF 已核验}
{\large\bfseries\color{Ink} RISC-V TEE 论文和规范发展快，survey 可减少分类断裂\par}\vspace{1.8mm}
\begin{columns}[T,totalwidth=\textwidth]
  \begin{column}{0.45\textwidth}
    \fcolorbox{Rule}{Paper}{%
  \begin{minipage}[t]{0.97\linewidth}
    \vspace{1.1mm}
    \hspace{1.4mm}\begin{minipage}{0.92\linewidth}
      \PanelTitle{内容说明}
      \scriptsize \begin{itemize}
  \item RISC-V TEE 论文和规范发展快，survey 可减少分类断裂
  \item 即使本地 PDF 已核验，机制细节仍要回到原始论文/spec
\end{itemize}
    \end{minipage}
    \vspace{1mm}
  \end{minipage}}
  \end{column}
  \begin{column}{0.49\textwidth}
    \fcolorbox{Rule}{Panel}{%
  \begin{minipage}[t]{0.97\linewidth}
    \vspace{1.1mm}
    \hspace{1.4mm}\begin{minipage}{0.92\linewidth}
      \PanelTitle{问题边界}
      \scriptsize {\tiny
\begin{tabularx}{\linewidth}{@{}YY@{}}
\toprule
\textbf{维度} & \textbf{说明} \\
\midrule
保护对象 & RISC-V CoVE / AP-TEE confidential VM \\
传统瓶颈 & 把 pre-CoVE enclave 与 AP-TEE/CoVE TVM lifecycle 区分开。 \\
本文入口 & RISC-V TEE 谱系辅助材料 \\
证据限制 & E2 survey / Background substrate \\
\bottomrule
\end{tabularx}
}
    \end{minipage}
    \vspace{1mm}
  \end{minipage}}
  \end{column}
\end{columns}
\vspace{1mm}
{\tiny\textcolor{Muted}{引用依据：引用依据为论文原文、官方规范或公开材料；本页结论按 E2 survey / Background substrate 的证据等级使用，不外推到未验证平台。}}
\end{frame}


\begin{frame}[t,shrink=6]{RISC-V TEE 谱系辅助材料：解决方案}
\DeckKicker{解决方案}
\SlideMeta{证据等级}{E2 survey / Background substrate}{来源状态：本地 PDF 已核验}
{\large\bfseries\color{Ink} 按 secure enclave、TEE mechanism、memory/I/O/attestation support 分类\par}\vspace{1.8mm}
\begin{columns}[T,totalwidth=\textwidth]
  \begin{column}{0.45\textwidth}
    \fcolorbox{Rule}{Paper}{%
  \begin{minipage}[t]{0.97\linewidth}
    \vspace{1.1mm}
    \hspace{1.4mm}\begin{minipage}{0.92\linewidth}
      \PanelTitle{内容说明}
      \scriptsize \begin{itemize}
  \item 按 secure enclave、TEE mechanism、memory/I/O/attestation support 分类
  \item 机制 结论 必须回引 CoVE paper 和 AP-TEE spec
\end{itemize}
    \end{minipage}
    \vspace{1mm}
  \end{minipage}}
  \end{column}
  \begin{column}{0.49\textwidth}
    \fcolorbox{Rule}{Panel}{%
  \begin{minipage}[t]{0.97\linewidth}
    \vspace{1.1mm}
    \hspace{1.4mm}\begin{minipage}{0.92\linewidth}
      \PanelTitle{核心设计拆解}
      \scriptsize \begin{itemize}
  \item 01. RISC-V enclave and TEE taxonomy
  \item 02. Pre-CoVE versus CoVE/AP-TEE classification
  \item 03. Survey bridge for representative systems
\end{itemize}
    \end{minipage}
    \vspace{1mm}
  \end{minipage}}
  \end{column}
\end{columns}
\vspace{1mm}
{\tiny\textcolor{Muted}{引用依据：引用依据为论文原文、官方规范或公开材料；本页结论按 E2 survey / Background substrate 的证据等级使用，不外推到未验证平台。}}
\end{frame}


\begin{frame}[t,shrink=6]{RISC-V TEE 谱系辅助材料：实验结果}
\DeckKicker{实验结果}
\SlideMeta{证据等级}{E2 survey / Background substrate}{来源状态：本地 PDF 已核验}
{\large\bfseries\color{Ink} 规范/Survey，无新实验；RISC-V TEE 谱系辅助材料 只支撑术语、taxonomy、接口语义或证据边界。\par}\vspace{1.8mm}
\begin{columns}[T,totalwidth=\textwidth]
  \begin{column}{0.45\textwidth}
    \fcolorbox{Rule}{Paper}{%
  \begin{minipage}[t]{0.97\linewidth}
    \vspace{1.1mm}
    \hspace{1.4mm}\begin{minipage}{0.92\linewidth}
      \PanelTitle{内容说明}
      \scriptsize \begin{itemize}
  \item 本页不展示独立系统实验，也不把二手归纳写成一手结果。
  \item 可支撑：Recent RISC-V TEE survey cross-checks the CoVE and AP-TEE lineage against older encla。
  \item 证据等级：E2 survey / Background substrate；来源状态：本地 PDF 已核验。
  \item 涉及具体平台、协议状态或数值时，转到原始论文、官方规范或厂商材料核验。
\end{itemize}
    \end{minipage}
    \vspace{1mm}
  \end{minipage}}
  \end{column}
  \begin{column}{0.49\textwidth}
    \fcolorbox{Rule}{Panel}{%
  \begin{minipage}[t]{0.97\linewidth}
    \vspace{1.1mm}
    \hspace{1.4mm}\begin{minipage}{0.92\linewidth}
      \PanelTitle{实验/证据边界}
      \scriptsize {\tiny
\begin{tabularx}{\linewidth}{@{}YYY@{}}
\toprule
\textbf{对象} & \textbf{可支撑} & \textbf{边界} \\
\midrule
数据/来源 & RISC-V TEE 谱系辅助材料 & 本地 PDF 已核验 \\
Baseline/对照 & 以论文或规范原文为准 & 不跨材料外推 \\
指标/对象 & 只使用当前来源可证明的信息 & E2 survey / Background substrate \\
使用边界 & 可进入本方向演进图 & 不能替代更强证据 \\
\bottomrule
\end{tabularx}
}
    \end{minipage}
    \vspace{1mm}
  \end{minipage}}
  \end{column}
\end{columns}
\vspace{1mm}
{\tiny\textcolor{Muted}{引用依据：引用依据为论文原文、官方规范或公开材料；本页结论按 E2 survey / Background substrate 的证据等级使用，不外推到未验证平台。}}
\end{frame}


\begin{frame}[t,shrink=6]{RISC-V TEE 谱系辅助材料：文章评价}
\DeckKicker{文章评价}
\SlideMeta{证据等级}{E2 survey / Background substrate}{来源状态：本地 PDF 已核验}
{\large\bfseries\color{Ink} 评价：RISC-V TEE 谱系辅助材料 适合做 taxonomy/规范入口；规范/Survey，无新实验，不能替代一手系统证据。\par}\vspace{1.8mm}
\begin{columns}[T,totalwidth=\textwidth]
  \begin{column}{0.45\textwidth}
    \fcolorbox{Rule}{Paper}{%
  \begin{minipage}[t]{0.97\linewidth}
    \vspace{1.1mm}
    \hspace{1.4mm}\begin{minipage}{0.92\linewidth}
      \PanelTitle{内容说明}
      \scriptsize \begin{itemize}
  \item 设计优势：有助于横向定位 Keystone/Penglai/SPEAR-V 与 CoVE/AP-TEE。
  \item 局限性：Survey 不能替代原始机制论文或官方规范，尤其不能承载性能或安全证明细节。
  \item 商业化潜力：适合产品路线图分类；风险在 survey 分类滞后和不同 threat model 被误合并。
  \item 规范/Survey，无新实验；具体平台结论转到一手材料核验。
\end{itemize}
    \end{minipage}
    \vspace{1mm}
  \end{minipage}}
  \end{column}
  \begin{column}{0.49\textwidth}
    \fcolorbox{Rule}{Panel}{%
  \begin{minipage}[t]{0.97\linewidth}
    \vspace{1.1mm}
    \hspace{1.4mm}\begin{minipage}{0.92\linewidth}
      \PanelTitle{文章评价}
      \scriptsize {\tiny
\begin{tabularx}{\linewidth}{@{}YY@{}}
\toprule
\textbf{维度} & \textbf{判断} \\
\midrule
设计优势 & 有助于横向定位 Keystone/Penglai/SPEAR-V 与 CoVE/AP-TEE。 \\
主要局限 & Survey 不能替代原始机制论文或官方规范，尤其不能承载性能或安全证明细节。 \\
商业落地 & 适合产品路线图分类；风险在 survey 分类滞后和不同 threat model 被误合并。 \\
讲稿定位 & 当前 SOTA/补充边界; 证据强度 E2 survey \\
\bottomrule
\end{tabularx}
}
    \end{minipage}
    \vspace{1mm}
  \end{minipage}}
  \end{column}
\end{columns}
\vspace{1mm}
{\tiny\textcolor{Muted}{引用依据：引用依据为论文原文、官方规范或公开材料；本页结论按 E2 survey / Background substrate 的证据等级使用，不外推到未验证平台。}}
\end{frame}


\begin{frame}[t,shrink=6]{RISC-V CoVE / AP-TEE confidential VM：方向总结}
\DeckKicker{方向总结}
\SlideMeta{证据等级}{方向综述}{来源状态：公开材料综合}
{\large\bfseries\color{Ink} RISC-V CoVE / AP-TEE confidential VM 的结论是：三篇主讲材料共同给出技术演进，但仍要用证据等级约束每个结论。\par}\vspace{1.8mm}
\begin{columns}[T,totalwidth=\textwidth]
  \begin{column}{0.45\textwidth}
    \fcolorbox{Rule}{Paper}{%
  \begin{minipage}[t]{0.97\linewidth}
    \vspace{1.1mm}
    \hspace{1.4mm}\begin{minipage}{0.92\linewidth}
      \PanelTitle{内容说明}
      \scriptsize \begin{itemize}
  \item RISC-V CoVE reference architecture：开创/基础入口，E3 / 基础证据 public preprint。
  \item RISC-V AP-TEE / CoVE TVM 规范草案：代表性改进，E3 / 草案/未批准规范。
  \item RISC-V TEE 谱系辅助材料：当前 SOTA/补充边界，E2 survey / Background substrate。
  \item 本方向结论：先用三篇主讲材料建立演进关系，再用证据等级控制机制、实验和商业化结论。
\end{itemize}
    \end{minipage}
    \vspace{1mm}
  \end{minipage}}
  \end{column}
  \begin{column}{0.49\textwidth}
    \fcolorbox{Rule}{Panel}{%
  \begin{minipage}[t]{0.97\linewidth}
    \vspace{1.1mm}
    \hspace{1.4mm}\begin{minipage}{0.92\linewidth}
      \PanelTitle{本方向方法对比}
      \scriptsize {\tiny
\begin{tabularx}{\linewidth}{@{}YYY@{}}
\toprule
\textbf{论文} & \textbf{定位} & \textbf{选择理由} \\
\midrule
RISC-V CoVE reference architecture & 开创/基础入口 & CoVE 2023 is the early mainline design material for RISC-V confidential V。 \\
RISC-V AP-TEE / CoVE TVM 规范草案 & 代表性改进 & AP-TEE draft is the current standards-track source for RISC-V confidentia。 \\
RISC-V TEE 谱系辅助材料 & 当前 SOTA/补充边界 & Recent RISC-V TEE survey cross-checks the CoVE and AP-TEE lineage against。 \\
\bottomrule
\end{tabularx}
}
    \end{minipage}
    \vspace{1mm}
  \end{minipage}}
  \end{column}
\end{columns}
\vspace{1mm}
{\tiny\textcolor{Muted}{引用依据：总结页综合三篇主讲材料的选择理由、评价字段和证据等级。}}
\end{frame}
